\section{Evaluation}

\label{sec:evaluation}

\subsection{Deployment Statistics}

Over 12 months of operation (January 2025 -- January 2026), ZDC accumulated:

\begin{table}[H]
\centering
\caption{Zoo Data Commons operational statistics.}
\label{tab:stats}
\begin{tabular}{@{}lr@{}}
\toprule
\textbf{Metric} & \textbf{Value} \\
\midrule
Total biodiversity records & 2,412,847 \\
Unique species represented & 14,237 \\
Contributing organizations & 847 \\
Geographic coverage (countries) & 94 \\
Total storage (IPFS) & 4.7 TB \\
Zoo Data Objects & 18,421 \\
Provenance records & 127,834 \\
Active validators & 312 \\
\bottomrule
\end{tabular}
\end{table}

\subsection{Data Quality Comparison}

We compared data quality scores between ZDC and three major centralized repositories using a standardized assessment framework applied to overlapping datasets:

\begin{table}[H]
\centering
\caption{Data quality comparison (mean scores, 0--1 scale).}
\label{tab:quality_comparison}
\begin{tabular}{@{}lcccc@{}}
\toprule
\textbf{Quality Dimension} & \textbf{ZDC} & \textbf{GBIF} & \textbf{iDigBio} & \textbf{iNaturalist} \\
\midrule
Metadata Completeness & 0.87 & 0.71 & 0.76 & 0.62 \\
Spatial Precision & 0.82 & 0.74 & 0.69 & 0.78 \\
Taxonomic Consistency & 0.91 & 0.83 & 0.87 & 0.71 \\
Temporal Resolution & 0.79 & 0.68 & 0.72 & 0.84 \\
Provenance Completeness & 0.94 & 0.42 & 0.51 & 0.33 \\
\midrule
\textbf{Overall} & \textbf{0.87} & 0.68 & 0.71 & 0.66 \\
\bottomrule
\end{tabular}
\end{table}

ZDC achieves 23.3\% higher overall quality than GBIF, with the largest advantage in provenance completeness (124\% higher), reflecting the architectural enforcement of provenance tracking.

\subsection{FAIR Assessment}

We evaluated FAIR compliance using the FAIRsFAIR assessment tool \citep{devaraju2021fairsfair}:

\begin{table}[H]
\centering
\caption{FAIR compliance scores (0--1 scale, assessed by FAIRsFAIR).}
\label{tab:fair}
\begin{tabular}{@{}lcccc@{}}
\toprule
\textbf{Principle} & \textbf{ZDC} & \textbf{GBIF} & \textbf{iDigBio} & \textbf{Dryad} \\
\midrule
Findable & 0.92 & 0.84 & 0.78 & 0.81 \\
Accessible & 0.88 & 0.91 & 0.83 & 0.87 \\
Interoperable & 0.85 & 0.79 & 0.74 & 0.72 \\
Reusable & 0.90 & 0.72 & 0.68 & 0.76 \\
\midrule
\textbf{Overall} & \textbf{0.89} & 0.82 & 0.76 & 0.79 \\
\bottomrule
\end{tabular}
\end{table}

ZDC scores highest on Findability (content-addressed identifiers) and Reusability (comprehensive provenance and licensing metadata), while GBIF maintains a slight edge on Accessibility due to its established API ecosystem.

\subsection{Access Control Efficacy}

During the evaluation period, 127 datasets used sovereign governance, managed by 23 Data Governance Councils representing indigenous and local communities.
Key metrics:

\begin{table}[H]
\centering
\caption{Sovereign access control statistics.}
\label{tab:sovereign}
\begin{tabular}{@{}lr@{}}
\toprule
\textbf{Metric} & \textbf{Value} \\
\midrule
Sovereign datasets & 127 \\
Data Governance Councils & 23 \\
Access requests processed & 412 \\
Requests approved & 287 (69.7\%) \\
Median response time & 4.2 days \\
Benefit-sharing agreements executed & 34 \\
TK labels applied & 89 \\
\bottomrule
\end{tabular}
\end{table}

The 69.7\% approval rate indicates that communities actively exercise governance authority rather than defaulting to open access, validating the need for the sovereign tier.

\subsection{Curation Market}

The curation market distributed 284,700 ZOO tokens during the evaluation period:

\begin{table}[H]
\centering
\caption{Curation market token distribution.}
\label{tab:market}
\begin{tabular}{@{}lrc@{}}
\toprule
\textbf{Activity} & \textbf{ZOO Distributed} & \textbf{Participants} \\
\midrule
Data Contribution & 142,300 & 847 \\
Validation & 67,400 & 312 \\
Annotation & 38,200 & 189 \\
Integration & 28,100 & 67 \\
Translation & 8,700 & 43 \\
\midrule
\textbf{Total} & 284,700 & -- \\
\bottomrule
\end{tabular}
\end{table}

The underserved region bonding curve directed 41\% of contribution rewards to data from Sub-Saharan Africa, Southeast Asia, and South America, regions historically underrepresented in global biodiversity databases.

\subsection{Cross-Institutional Reuse}

ZDC facilitated 312 cross-institutional data reuse events over 12 months, compared to an estimated 89 events that would have occurred through traditional channels (based on pre-deployment reuse rates for the same research groups).
These reuse events supported 47 peer-reviewed publications, 12 conservation management plans, and 8 environmental impact assessments.

% ============================================================
% 9. Discussion
% ============================================================
